‫\فصل{مقدمه}
‫
‫نخستین فصل این پایان‌نامه به معرفی مسئله، بیان اهمیت و ضرورت موضوع، تبیین چالش‌های موجود در معماری‌های رایج، اهداف پژوهش و در نهایت معرفی ساختار کلی پایان‌نامه می‌پردازد. در این فصل، تصویری جامع از انگیزه‌های انجام این تحقیق و مسیر پیش رو ارائه می‌گردد.
‫
‫\قسمت{بیان مسئله}
‫
‫در سال‌های اخیر، ظهور و تکامل مدل‌های زبانی بزرگ\پانویس{Large Language Models (LLMs)} تحولی بنیادین و پارادایم‌شیفت چشمگیری در نحوه تعامل انسان با ماشین، پردازش زبان طبیعی و روش‌های دسترسی به اطلاعات ایجاد کرده است. با این وجود، پارادایم سنتیِ بازیابی اطلاعات و جستجو در وب همچنان با چالش‌ها و کاستی‌های جدی مواجه است. موتورهای جستجوی سنتی، عمدتاً بر اساس الگوریتم‌های تطبیق کلیدواژه\پانویس{Keyword-matching} و تحلیل پیوندها عمل می‌کنند. این سیستم‌ها به جای ارائه پاسخ نهایی و تحلیلی به کاربر، صرفاً فهرستی از پیوندها (لینک‌ها) را برمی‌گردانند که بارِ شناختیِ مطالعه، تجمیع و استخراج پاسخ نهایی را به طور کامل بر دوش کاربر قرار می‌دهد.
‫
‫از سوی دیگر، مدل‌های زبانی بزرگ با وجود توانمندی شگرف در درک عمیق بافتار و تولید زبان طبیعیِ روان، دارای محدودیت‌های ذاتی و ساختاریِ مهمی هستند. نخستین محدودیت، اتکای این مدل‌ها به «دانش ایستا»\پانویس{Static Knowledge} است؛ به این معنا که حافظه پارامتریک\پانویس{Parametric Memory} آن‌ها محدود به داده‌هایی است که تا زمان پایان فاز آموزش\پانویس{Training Cut-off} دریافت کرده‌اند. به همین دلیل، این مدل‌ها در پاسخ‌گویی به پرسش‌هایی که نیازمند آگاهی از رویدادهای لحظه‌ای، اخبار روز و داده‌های بلادرنگ هستند، کاملاً ناتوان عمل می‌کنند. 
‫
‫محدودیت دوم و به مراتب خطرناک‌تر، پدیده‌ای به نام «توهم»\پانویس{Hallucination} است. مدل‌های زبانی به گونه‌ای طراحی شده‌اند که همواره توالی کلمات بعدی را با بیشترین احتمال تولید کنند؛ بنابراین در صورت عدم آگاهی از فکت‌های واقعی پیرامون یک موضوع، ممکن است پاسخ‌هایی کاملاً ساختگی، نامعتبر اما با لحنی به شدت متقاعدکننده و ساختاریافته تولید نمایند. 
‫
‫بر این اساس، مسئله اصلی پژوهش حاضر، غلبه بر محدودیت‌های متقابل هر دو سیستم (یعنی رویکرد منفعل موتورهای جستجوی سنتی و توهم‌زایی مدل‌های زبانی منفرد) از طریق طراحی، پیاده‌سازی و بررسی یک سامانه پاسخ‌گویی هوشمند است. سامانه‌ای عامل‌گونه که بتواند با تحلیل دقیق نیت کاربر و درک تاریخچه تعاملات (گفتگوها)، پرسش‌های جستجوی بهینه‌ای تولید کرده، فضای وب را در لحظه کاوش کند و با بهره‌گیری از تکنیک‌های پیشرفته پردازش زبان و بازرتبه‌بندی\پانویس{Reranking}، پاسخی تحلیلی، به‌روز و کاملاً مستند (همراه با ارجاع به منابع) ارائه دهد.
‫
‫\قسمت{اهمیت و ضرورت موضوع}
‫
‫با تغییر الگوی مصرف اطلاعات در عصر دیجیتال، سطح انتظارات کاربران از سیستم‌های بازیابی اطلاعات دگرگون شده است. کاربران امروزی انتظار دارند به جای جستجوی دستی در میان انبوهی از صفحات وب، پاسخ‌های مستقیم، دقیق، به‌روز و قابل استناد را در قالب یک گفتگوی محاوره‌ای و روان دریافت کنند. ترکیب موتورهای جستجو با هوش مصنوعی مولد\پانویس{Generative AI}، راهکاری است که معماری‌های پیشگامی نظیر پرپلکسیتی\پانویس{Perplexity}، بینگ کوپایلت\پانویس{Bing Copilot} و جستجوی هوشمند گوگل\پانویس{Google AI Search} به سوی آن حرکت کرده‌اند و اهمیت استراتژیک و تجاری این حوزه را در سطح جهانی به اثبات رسانده‌اند.
‫
‫ضرورت بنیادین انجام این پژوهش در این است که استفاده‌ی صرف و بدون پشتوانه از مدل‌های زبانی در حوزه‌های حساس (نظیر حوزه‌های پزشکی، حقوقی، علمی و خبری) که نیازمند دقت بالا و فکت‌های اثبات‌شده هستند، به دلیل خطرات ناشی از توهم مدل، غیرقابل اعتماد و بعضاً آسیب‌زاست. متصل کردن مدل زبانی به جریان اطلاعات پویای وب که اصطلاحاً اتصال به زمینه‌ی وب\پانویس{Web Grounding} نامیده می‌شود، نه تنها معضل اطلاعات منقضی‌شده را برطرف می‌سازد، بلکه با فراهم آوردن امکان ارجاع‌دهی دقیق به منابع و مراجع معتبر\پانویس{Citations}، شفافیت و اعتماد کاربر به سیستم را به شکل چشمگیری افزایش می‌دهد.
‫
‫\قسمت{چالش‌های معماری استاندارد RAG در فضای وب}
‫
‫در ادبیات موضوع فعلی، معماری «تولید افزوده با بازیابی» یا رگ\پانویس{Retrieval-Augmented Generation (RAG)} به عنوان راهکاری استاندارد و اثبات‌شده برای تزریق دانش خارجی و غیرپارامتریک به مدل‌های زبانی شناخته می‌شود. با این حال، خط لوله‌ی\پانویس{Pipeline} RAG سنتی عمدتاً برای کار بر روی اسناد ثابت، ساختاریافته و آفلاین (مانند پایگاه داده‌ای از فایل‌های سازمانی) طراحی شده است و در مواجهه با فضای وب با چالش‌های فنی و معنایی اساسی روبرو است.
‫فضای وب، محیطی به شدت پویا، پرنویز، ناهمگن و غیرقابل پیش‌بینی است. در RAG سنتی، اسناد پیش از استفاده، طی یک فرآیند آفلاین پردازش، پاک‌سازی و قطعه‌بندی\پانویس{Chunking} شده و در یک پایگاه داده برداری\پانویس{Vector Database} ذخیره می‌گردند. اما در فرآیند وب‌گردی در لحظه\پانویس{Real-time Web Browsing}، زمان کافی برای این پیش‌پردازشِ عمیق وجود ندارد. سیستم باید بتواند در کسری از ثانیه، محتوای صفحات وبِ بازیابی‌شده را که حاوی حجم زیادی از تبلیغات، کدهای مخرب، ساختارهای پیچیده HTML و نویزهای بصری هستند، دریافت و پالایش کند. 
‫
‫همچنین در RAG سنتی، معمولاً پرسش اولیه کاربر مستقیماً برای جستجوی معنایی استفاده می‌شود؛ اما در وب، پرسش اولیه کاربر ممکن است مبهم یا چندوجهی باشد و برای یافتن پاسخ جامع، نیازمند تجزیه به چندین پرس‌وجوی فرعی و مستقل\پانویس{Multi-query Generation} باشد؛ قابلیتی که در سیستم‌های جستجوی معناییِ پایه تعبیه نشده و نیازمند یک ماژول تصمیم‌گیر (عامل هوشمند) برای مدیریت بسط پرس‌وجوهاست.
‫
‫\قسمت{اهداف پژوهش}
‫
‫هدف اصلی این پژوهش، طراحی، توسعه و ارزیابی یک موتور جستجوی هوشمند و گفتگومحور است که توانایی استخراج بلادرنگ اطلاعات از وب و ترکیب آن با قدرت استدلال مدل‌های زبانی را داشته باشد. در راستای تحقق این هدف کلان، اهداف خرد زیر دنبال می‌شوند:
‫
‫\شروع{فقرات}
‫\فقره
‫مطالعه، بررسی و تحلیل معماری سامانه‌های نوین جستجوی هوشمند مبتنی بر مدل‌های زبانی بزرگ.
‫\فقره
‫طراحی یک سامانه ترکیبیِ دوحالته با قابلیت تشخیص خودکار و سوئیچ بین «پاسخ‌گویی عادی (مبتنی بر دانش پارامتریک)» و «پاسخ‌گویی مبتنی بر بازیابی از وب».
‫\فقره
‫پیاده‌سازی مکانیزم آگاهی از بافتار\پانویس{Context-awareness} با استفاده از تزریق تاریخچه گفتگوها جهت تولید پاسخ‌های پیوسته، منطقی و چندنوبتی\پانویس{Multi-turn}.
‫\فقره
‫بررسی و پیاده‌سازی مکانیزم‌های پیشرفته‌ی بازنمره‌دهی و بازرتبه‌بندی پرس‌وجوها و متون استخراج‌شده از وب برای یافتن مرتبط‌ترین قطعات متنی و کاهش نویز.
‫\فقره
‫طراحی یک چارچوب ارزیابی جامع برای سنجش عملکرد و کیفیت خروجی سامانه پیشنهادی روی مجموعه‌ای از پرسش‌های چالشی، مبهم و نیازمند اطلاعات لحظه‌ای.
‫\پایان{فقرات}
‫
‫\قسمت{پرسش‌های پژوهش}
‫
‫این پژوهش در پی یافتن پاسخ علمی، مستدل و قابل اندازه‌گیری برای پرسش‌های اساسی زیر است:
‫
‫\شروع{فقرات}
‫\فقره
‫آیا ادغام جستجوی وبِ بلادرنگ با مدل‌های زبانی بزرگ، موجب کاهش معنادار پدیده توهم و بهبود دقت فاکچوال\پانویس{Factual Accuracy} در پاسخ‌گویی به پرسش‌های دانش‌محور می‌شود؟
‫\فقره
‫آیا به‌کارگیری تکنیک‌ها و مدل‌های بازرتبه‌بندی\پانویس{Reranker} بر روی نتایج خامِ استخراج‌شده از صفحات وب، کیفیت، انسجام و ارتباط معنایی پاسخ نهایی ارائه‌شده به کاربر را ارتقا می‌دهد؟
‫\فقره
‫در یک محیط تعاملی و چندنوبتی، تزریق تاریخچه گفتگوها به ماژول تولید پرس‌وجو، تا چه میزان در درک بهتر نیت پنهان کاربر و تولید کلیدواژه‌های جستجوی دقیق‌ترِ مؤثر است؟
‫\پایان{فقرات}

‫\قسمت{نوآوری و سهم کار}
‫
‫سهم اصلی و نوآوری این پژوهش، فراتر رفتن از یک رابط برنامه‌نویسیِ\پانویس{API} ساده و توسعه یک سیستم یکپارچه، ماژولار و عامل‌گونه\پانویس{Agentic} است که به صورت خودمختار توانایی مدیریت جریان اطلاعات را داراست. دستاوردهای کلیدی و نوآوری‌های این پایان‌نامه عبارتند از:
‫
‫\شروع{فقرات}
‫\فقره
‫طراحی و پیاده‌سازی یک سامانه‌ی پردازشی هوشمند با قابلیت تشخیص خودکار نیاز به جستجوی خارجی (کاهش سربار سیستم با طراحی دوحالته).
‫\فقره
‫توسعه یک ماژول تخصصی برای تولید و بسط پرس‌وجو\پانویس{Query Expansion}، که قادر است با تحلیل ترکیب سؤال فعلی کاربر و بافتارِ تاریخچه چت، چندین پرس‌وجوی مستقل و بهینه برای موتور جستجوی وب تولید نماید.
‫\فقره
‫توسعه الگوریتمی برای تجمیع اطلاعات پراکنده از منابع مختلف وب، پالایش نویزهای HTML، و تولید یک پاسخ نهایی جامع همراه با سیستم ارجاع‌دهی درون‌متنی\پانویس{In-text Citation}.
‫\فقره
‫طراحی و اجرای یک خط لوله ارزیابی علمی برای مقایسه عملکرد سیستم در حالت پایه\پانویس{Basic Web Retrieval} در برابر حالت مجهز به مدل‌های بازرتبه‌بندی.
‫\پایان{فقرات}
‫
‫\قسمت{روش کلی انجام کار}
‫
‫روند کلی انجام این پژوهش مبتنی بر توسعه یک خط لوله (Pipeline) پردازشیِ چندمرحله‌ای و هوشمند است. در این فرآیند، ابتدا درخواست کاربر دریافت شده و همراه با تاریخچه گفتگوهای پیشین به ماژول تحلیل نیت ارسال می‌گردد. سیستم با بهره‌گیری از یک مدل زبانی پایه، نیاز اطلاعاتی کاربر را تحلیل کرده و در صورت تشخیص نیاز به داده‌های به‌روز، آن را به چند پرس‌وجوی (Query) بهینه و مجزا می‌شکند. 
‫
‫در گام بعدی، این پرس‌وجوها از طریق رابط موتور جستجو به شبکه وب ارسال شده و صفحات مرتبط شناسایی و بارگیری (Scrape) می‌شوند. محتوای متنی این صفحات استخراج شده و به قطعات کوچکتر تقسیم می‌گردد. سپس به کمک یک مدل بازرتبه‌بند مبتنی بر شبکه‌های عصبی متقاطع\پانویس{Cross-encoder Reranker}، این قطعات متنی بر اساس بیشترین میزان هم‌پوشانی معنایی با نیاز کاربر، فیلتر و مرتب‌سازی می‌شوند. 
‫
‫در نهایت، متونِ پالایش‌شده‌ی برتر به عنوان «بافتارِ دانش» به دستورالعملِ\پانویس{Prompt} مدل زبانی بزرگ تزریق می‌شوند تا مدل بتواند پاسخی نهایی، دقیق، مستدل و ارجاع‌دار تولید کرده و به کاربر نمایش دهد. در فاز پایانی پژوهش نیز، کیفیت خروجی‌های این معماری با استفاده از متریک‌های ارزیابی کیفیِ رایج در حوزه پردازش زبان طبیعی سنجیده و تحلیل خواهند شد.
‫
‫\قسمت{ساختار پایان‌نامه}
‫
‫[این بخش پس از نهایی شدن نگارش فصول بعدی پایان‌نامه تکمیل خواهد شد. در این قسمت طی یک یا دو پاراگراف، محتوای فصول دوم (مروری بر ادبیات موضوع)، سوم (روش‌شناسی و معماری سیستم پیشنهادی)، چهارم (پیاده‌سازی، آزمایش‌ها و ارزیابی نتایج) و پنجم (نتیجه‌گیری و کارهای آتی) به اختصار شرح داده خواهد شد.]
‫